% =============================================================
%  EEE 2188 – Electronics Sessional
%  Experiment 01: Familiarization of Equipment
%  Rajshahi University of Engineering & Technology (RUET)
% =============================================================

\documentclass[12pt, a4paper, twoside]{report}

% ---------------------------------------------------------------
% Core packages
% ---------------------------------------------------------------
\usepackage[T1]{fontenc}
\usepackage[utf8]{inputenc}
\usepackage{mathptmx}           % Times New Roman for text + math
\usepackage{microtype}          % Better text justification

% ---------------------------------------------------------------
% Page geometry
% ---------------------------------------------------------------
\usepackage[
  a4paper,
  top=2.5cm, bottom=2.5cm,
  inner=3.0cm, outer=2.5cm,
  headheight=15pt
]{geometry}

% ---------------------------------------------------------------
% Headers / Footers
% ---------------------------------------------------------------
\usepackage{fancyhdr}
\usepackage{ifthen}
\usepackage{lastpage}

% Custom marks so headers update per section
\pagestyle{fancy}
\fancyhf{}

% Even pages (left page): Course Code  Experiment  No.
% Odd  pages (right page): Course Title  First Section on page
\fancyhead[LE]{\small\textit{\CourseCode\ \ Experiment\ \ExperimentNo}}
\fancyhead[RO]{\small\textit{\CourseTitle\ \ \rightmark}}
\fancyfoot[C]{\small\thepage}
\renewcommand{\headrulewidth}{0.4pt}

% Keep \rightmark updated from \section commands
\renewcommand{\sectionmark}[1]{\markright{#1}}

% Plain style for chapter-opening pages
\fancypagestyle{plain}{%
  \fancyhf{}%
  \fancyfoot[C]{\small\thepage}%
  \renewcommand{\headrulewidth}{0pt}%
}

% ---------------------------------------------------------------
% Figures, tables, captions
% ---------------------------------------------------------------
\usepackage{graphicx}
\usepackage{float}
\usepackage[
  format=plain,
  labelfont=bf,
  font=small,
  justification=centering
]{caption}

% Chapter-based figure/table numbering: Fig. 1.1, Fig. 1.2 …
\usepackage{chngcntr}
\counterwithin{figure}{chapter}
\counterwithin{table}{chapter}

% ---------------------------------------------------------------
% Colours, boxes
% ---------------------------------------------------------------
\usepackage[dvipsnames]{xcolor}
\usepackage{tcolorbox}
\tcbuselibrary{skins, breakable}

% A subtle definition/term box
\tcbset{
  termbox/.style={
    enhanced,
    colback=gray!8,
    colframe=gray!50,
    boxrule=0.6pt,
    arc=3pt,
    left=8pt, right=8pt, top=6pt, bottom=6pt,
    breakable
  }
}

% ---------------------------------------------------------------
% Hyperlinks
% ---------------------------------------------------------------
\usepackage[
  colorlinks=true,
  linkcolor=black,
  citecolor=black,
  urlcolor=blue,
  pdfborder={0 0 0}
]{hyperref}

% ---------------------------------------------------------------
% Bibliography – IEEE style
% ---------------------------------------------------------------
\usepackage[numbers, sort&compress]{natbib}
\bibliographystyle{ieeetr}

% ---------------------------------------------------------------
% Spacing / Paragraph
% ---------------------------------------------------------------
\usepackage{setspace}
\setstretch{1.25}
\setlength{\parskip}{6pt}
\setlength{\parindent}{0pt}

% ---------------------------------------------------------------
% Section heading styles (no bold font weight explosion)
% ---------------------------------------------------------------
\usepackage{titlesec}
\titleformat{\chapter}[display]
  {\normalfont\large\bfseries\centering}
  {\MakeUppercase{Experiment No.~\thechapter}}{6pt}
  {\Large\MakeUppercase}
\titlespacing*{\chapter}{0pt}{10pt}{20pt}

\titleformat{\section}
  {\normalfont\normalsize\bfseries}{\thesection}{1em}{}
\titleformat{\subsection}
  {\normalfont\normalsize\itshape}{\thesubsection}{1em}{}

% ---------------------------------------------------------------
% Document-level metadata (edit here for new experiments)
% ---------------------------------------------------------------
\newcommand{\CourseCode}{EEE~2188}
\newcommand{\CourseTitle}{Electronics Sessional}
\newcommand{\ExperimentNo}{01}
\newcommand{\ExperimentName}{Familiarization of Different Equipment\\[4pt]Involved with Electronics Sessional}
\newcommand{\Department}{Department of Electrical \& Electronic Engineering}
\newcommand{\Institution}{Rajshahi University of Engineering \& Technology}

% ---------------------------------------------------------------
% Placeholder image macro
% Usage: \imgplaceholder{caption text}{label}
% ---------------------------------------------------------------
\newcommand{\imgplaceholder}[2]{%
  \begin{figure}[H]
    \centering
    \fcolorbox{gray!60}{gray!12}{%
      \parbox[c][5cm][c]{0.65\textwidth}{%
        \centering
        \textcolor{gray!80}{\Large [Image Placeholder]}\\[6pt]
        \textcolor{gray}{\small Insert figure here}
      }%
    }
    \caption{#1}
    \label{fig:#2}
  \end{figure}
}

% =============================================================
\begin{document}
% =============================================================

% ---------------------------------------------------------------
% Title page
% ---------------------------------------------------------------
\begin{titlepage}
  \centering
  \vspace*{1.5cm}
  {\Large\bfseries \CourseTitle \par}
  \vspace{0.4cm}
  {\large\bfseries (\CourseCode) \par}
  \vspace{1.2cm}
  \rule{\linewidth}{0.6pt}\\[8pt]
  {\LARGE\bfseries Experiment No.~\ExperimentNo \par}
  \vspace{0.6cm}
  {\Large \ExperimentName \par}
  \rule{\linewidth}{0.6pt}
  \vfill
  {\large \Department \par}
  \vspace{0.3cm}
  {\large \Institution \par}
  \vspace{1.5cm}
\end{titlepage}

% ---------------------------------------------------------------
% Chapter: Experiment 01
% ---------------------------------------------------------------
\chapter[\ExperimentNo]{\ExperimentName}
\markright{Objectives}   % initialise right-mark

% =============================================================
\section{Objectives}
\markright{Objectives}
% =============================================================

The objectives of this experiment are to learn about the commonly used equipment in the
Electronics Sessional, to observe the various parameters of an electric circuit, to identify
the different instruments ordinarily employed for measurement, and to understand the use
of these instruments and components in experimental work.

% =============================================================
\section{Theory}
\markright{Theory}
% =============================================================

Electronics is a scientific and engineering discipline that studies and applies the
principles of physics in order to design, create, and operate devices capable of
manipulating electrons and other electrically charged particles \cite{sedra2020}.
As a subfield of physics and electrical engineering, it exploits active devices---such as
transistors, diodes, and integrated circuits---to control and amplify the flow of electric
current, and to convert signals from one form to another: for instance, from alternating
current (AC) to direct current (DC), or from analog signals to digital ones.

The evolution of electronics over the last century has been driven by a relentless
reduction in component size and cost alongside dramatic improvements in speed, efficiency,
and reliability. Beginning with vacuum tubes in the early twentieth century, the field
transitioned first to discrete semiconductor devices and then to the densely integrated
circuits that now underpin everything from household appliances to satellite communication
systems. Each generation of components is governed by the same fundamental physical
laws---Ohm's law, Kirchhoff's laws, and Maxwell's equations---yet their practical
realization has changed enormously. Understanding the individual components at the heart
of any electronic circuit is therefore the essential first step for any engineering
student.

In a typical electronics laboratory, a circuit under investigation consists of passive
components (resistors, capacitors) that store or dissipate energy, and active components
(diodes, transistors, operational amplifiers) that can amplify or switch signals
\cite{boylestad2015}. Measurement instruments such as the multimeter and the oscilloscope
allow the engineer to probe voltages, currents, and waveforms at any node of the circuit,
while prototyping boards---the breadboard and the veroboard---provide flexible platforms
for assembling circuits without committing to a permanent printed-circuit-board (PCB)
layout. The soldering iron, finally, is the hand tool by which components are permanently
bonded to a PCB once a design has been validated. Together, these components and
instruments form the core toolkit of the electronics practitioner, and familiarity with
each of them is the primary goal of this experiment.

% =============================================================
\section{Major Components and Equipment}
\markright{Major Components and Equipment}
% =============================================================

% -----------------------------------------------------------
\subsection{Resistor}
\markright{Resistor}
% -----------------------------------------------------------

\textbf{Definition.}\quad
A resistor is a passive, two-terminal electronic component that implements electrical
resistance as a circuit element \cite{wiki:resistor}. It obeys Ohm's Law ($V = IR$),
meaning the voltage drop across it is directly proportional to the current flowing through
it. Resistance is measured in ohms ($\Omega$), and resistors values are most commonly
indicated by a system of color-coded bands printed on their body.

\imgplaceholder{Resistor and color-band coding chart}{resistor}

\textbf{Specifications.}\quad
The primary specifications of a resistor are its resistance value (ranging from fractions
of an ohm up to several gigaohms), its tolerance (the permissible deviation from the
nominal value, commonly $\pm1\%$, $\pm5\%$, or $\pm10\%$), and its power rating (the
maximum heat it can safely dissipate, typically $\frac{1}{8}$\,W to $2$\,W for through-hole
types and up to several hundred watts for industrial power resistors) \cite{electronics_notes:resistor}.
The temperature coefficient of resistance (TCR), expressed in parts per million per degree
Celsius (ppm/$^\circ$C), quantifies how much the resistance drifts with temperature;
carbon-film resistors exhibit TCR values of $\pm100$ to $\pm500$\,ppm/$^\circ$C, while
metal-film types achieve as low as $\pm25$\,ppm/$^\circ$C, and precision metal-foil
resistors can reach $\pm0.2$\,ppm/$^\circ$C \cite{electronics_notes:resistor}.
Resistors are available in through-hole (axial lead) and surface-mount (SMD) packages;
SMD types are labeled using a three- or four-digit numeric code rather than color bands.

\textbf{Purpose of Use.}\quad
The essential purpose of a resistor is to limit or regulate the flow of electric current
in a circuit. Without current-limiting resistors, sensitive components such as LEDs and
transistors would draw destructive levels of current. Beyond simple current limitation,
resistors are used to divide voltage (voltage dividers), set the gain of amplifiers, bias
the operating point of transistors, terminate transmission lines to prevent signal
reflection, and provide pull-up or pull-down functionality on digital logic lines.

\textbf{Applications.}\quad
Resistors are ubiquitous in virtually every electronic system. They appear in signal
conditioning networks in communication devices, load regulators in power supplies, sensor
networks for heat, light, and force sensing (thermistors, LDRs, strain gauges), biasing
circuits for BJTs and MOSFETs, and protection networks at the inputs and outputs of
integrated circuits \cite{boylestad2015}. In power electronics, large wirewound or
grid-type resistors are used for dynamic braking in motor drives and for load-bank testing
of generators.

% -----------------------------------------------------------
\subsection{Multimeter}
\markright{Multimeter}
% -----------------------------------------------------------

\textbf{Definition.}\quad
A multimeter---also known as a volt-ohm meter (VOM)---is a multi-function electronic
measuring instrument that combines the capabilities of a voltmeter, an ammeter, and an
ohmmeter in a single handheld or bench-mounted device \cite{wiki:multimeter}. Modern
digital multimeters (DMMs) display readings on a numeric LCD, whereas the earlier analog
multimeter uses a moving-coil mechanism with a pointer deflecting over a graduated scale.

\imgplaceholder{Digital multimeter (front view with controls)}{multimeter}

\textbf{Specifications.}\quad
A typical handheld DMM measures DC voltage in ranges from 200\,mV to 600\,V, AC voltage
up to 600\,V, DC current from 200\,$\mu$A to 10\,A, resistance from 200\,$\Omega$ to
$2\,\text{M}\Omega$, and sometimes capacitance and frequency as well \cite{electronicsclub:multimeter}.
Standard portable DMMs achieve a DC voltage accuracy of $\pm0.5\%$ of reading, while
laboratory-grade bench instruments can attain accuracies better than $\pm0.01\%$; analog
multimeters typically carry a less precise rating of $\pm3\%$ full-scale deflection
\cite{wiki:multimeter}. Input impedance on most DMMs is fixed at $10\,\text{M}\Omega$,
high enough to avoid loading down the circuit under test.

\textbf{Purpose of Use.}\quad
The multimeter is used to verify the operating voltages and currents at nodes of a
circuit, to check component values (resistance, capacitance), to test diodes and
transistors, and to measure continuity---that is, whether an electrical connection is
sound. When used as a voltmeter, it is connected in parallel with the element under test;
when used as an ammeter, it must be inserted in series, typically using a low-value shunt
resistor and a dedicated 10\,A terminal.

\textbf{Applications.}\quad
Multimeters are indispensable in electronics laboratories for circuit debugging,
component testing, and quality verification. They are equally essential in field service
for troubleshooting household wiring, automotive electrical systems, industrial motor
drives, and battery-powered consumer devices \cite{wiki:multimeter}. Specialized
multimeters---sometimes called \textit{clamp meters}---measure AC current through a wire
without breaking the circuit, which is invaluable for high-current industrial diagnostics.

% -----------------------------------------------------------
\subsection{Project Board / Breadboard}
\markright{Project Board / Breadboard}
% -----------------------------------------------------------

\textbf{Definition.}\quad
A breadboard (or solderless breadboard) is a reusable construction base used to build
and test prototype electronic circuits without requiring any soldering \cite{wiki:breadboard}.
It consists of a perforated plastic block housing an array of spring-clip contacts beneath
the holes. Components are inserted into the holes and held by friction, while internal
copper strips connect groups of holes to form electrical nodes.

\imgplaceholder{Solderless breadboard showing internal connection pattern}{breadboard}

\textbf{Specifications.}\quad
The holes on a standard breadboard are spaced 0.1\,inches (2.54\,mm) apart, matching the
pin pitch of common dual in-line package (DIP) integrated circuits. A full-size breadboard
typically offers 830 tie-points arranged in two terminal strips of 63 rows each and two
pairs of power rails running along the edges \cite{wiki:breadboard}. Practical breadboards
exhibit parasitic capacitance of a few picofarads and contact resistance of a few ohms;
these parasitics limit reliable operation to frequencies below approximately 10\,MHz
\cite{analog:breadboard}.

\textbf{Purpose of Use.}\quad
The breadboard allows engineers and students to assemble and modify circuits quickly,
test different component configurations, and verify circuit behavior---all without the
commitment and expense of a permanent soldered PCB. Its reusable nature makes it
particularly suitable for iterative design and experimentation: a component placed
incorrectly can be moved to another position in seconds.

\textbf{Applications.}\quad
Breadboards are used throughout electronics education, prototyping, and research. They
are the standard platform for testing analog filter designs, digital logic circuits,
microcontroller-based systems (such as Arduino and Raspberry Pi), and RF circuits at
moderate frequencies \cite{analog:breadboard}. Once a design is validated on a
breadboard, it is transferred to a permanent PCB or, for intermediate permanence,
to a veroboard.

% -----------------------------------------------------------
\subsection{Capacitor}
\markright{Capacitor}
% -----------------------------------------------------------

\textbf{Definition.}\quad
A capacitor is a passive two-terminal electronic component that stores and releases
electrical energy in an electric field \cite{sedra2020}. It consists of two conductive
plates separated by an insulating material called the dielectric. The ability to store
charge is quantified by its capacitance $C$, defined as $C = Q/V$, where $Q$ is the
stored charge and $V$ is the voltage across the plates.

\imgplaceholder{Electrolytic capacitor (radial type, 100\,$\mu$F / 400\,V)}{capacitor}

\textbf{Specifications.}\quad
Capacitance values span an enormous range---from a few femtofarads (fF) for small ceramic
chip capacitors to tens of thousands of microfarads ($\mu$F) for large electrolytic
types. Key specifications include capacitance value, voltage rating (the maximum DC
voltage that can be sustained without breakdown), temperature coefficient (how capacitance
changes with temperature), equivalent series resistance (ESR), and leakage current.
Common dielectric materials include ceramic (stable, low loss, suitable for high-frequency
bypass), electrolytic oxide film (high capacitance density, polarized), tantalum oxide
(compact, polarized, low ESR), and polyester film (non-polarized, general-purpose)
\cite{boylestad2015}.

\textbf{Purpose of Use.}\quad
Capacitors serve several distinct roles depending on the circuit context. They store
energy and release it when needed (energy storage), block DC while passing AC signals
(coupling and DC blocking), smooth out ripple in power supply outputs (filtering),
tune resonant circuits (in combination with inductors), and create time delays in RC
timing circuits.

\textbf{Applications.}\quad
Capacitors are found in every electronic sub-system. Power supply filter capacitors
reduce ripple on DC rails; bypass capacitors suppress high-frequency noise on power pins
of ICs; coupling capacitors transfer audio or RF signals between stages without
passing the DC bias; and timing capacitors set oscillation frequencies in clock circuits
and astable multivibrators \cite{horowitz2015}.

% -----------------------------------------------------------
\subsection{Diode}
\markright{Diode}
% -----------------------------------------------------------

\textbf{Definition.}\quad
A diode is a two-terminal semiconductor device that allows electric current to flow
predominantly in one direction, thereby functioning as a one-way valve for current
\cite{boylestad2015}. It is formed by a junction of p-type and n-type semiconductor
material (a p--n junction). When the anode is positive with respect to the cathode
(forward bias), the diode conducts with a typical forward voltage drop of about 0.6--0.7\,V
for silicon; in the reverse direction (reverse bias), it blocks current until the reverse
breakdown voltage is exceeded.

\imgplaceholder{Rectifier diode (1N4003) showing physical package and schematic symbol}{diode}

\textbf{Specifications.}\quad
Key diode specifications are the forward voltage drop ($V_F$), the maximum forward
current ($I_F$), the peak inverse voltage (PIV) or reverse breakdown voltage, and the
reverse saturation current ($I_S$). A standard rectifier diode such as the 1N4007 has
$V_F \approx 1.1$\,V at 1\,A, PIV of 1000\,V, and an average forward current rating
of 1\,A. Fast-recovery diodes and Schottky diodes offer lower $V_F$ (0.2--0.4\,V for
Schottky) and faster switching times for high-frequency applications \cite{sedra2020}.

\textbf{Purpose of Use.}\quad
The primary use of a diode is rectification---converting AC to DC by allowing only
one polarity of current to pass. Diodes are also used for voltage clamping (protecting
circuits from voltage spikes), signal demodulation in radio receivers, logic level
shifting, and, in the case of light-emitting diodes, the generation of visible or infrared
light.

\textbf{Applications.}\quad
Diodes appear in AC--DC power supply rectifier bridges, battery chargers (to prevent
reverse current), signal detectors in AM radio receivers, freewheeling or flyback
circuits that protect switching transistors from inductive voltage spikes, and overvoltage
protection circuits on the inputs of logic ICs \cite{floyd2012}.

% -----------------------------------------------------------
\subsection{Zener Diode}
\markright{Zener Diode}
% -----------------------------------------------------------

\textbf{Definition.}\quad
A Zener diode is a special type of silicon p--n junction diode designed to operate
reliably in the reverse-bias breakdown region \cite{boylestad2015}. Unlike a conventional
rectifier diode, which is destroyed by sustained reverse breakdown, the Zener diode is
deliberately constructed so that once the applied reverse voltage reaches its characteristic
\textit{Zener voltage} ($V_Z$), it conducts in the reverse direction while maintaining
that voltage nearly constant over a wide range of current. This phenomenon, first described
by American physicist Dr.\ Clarence Melvin Zener, makes the device invaluable as a voltage
reference.

\imgplaceholder{Zener diode in TO-92 package with schematic symbol}{zener}

\textbf{Specifications.}\quad
The most important specification of a Zener diode is its Zener (breakdown) voltage $V_Z$,
which is available in standard values ranging from about 2.4\,V to 200\,V. Additional
parameters include the power dissipation rating ($P_Z$), the dynamic impedance ($Z_Z$,
typically a few ohms, indicating how stable the voltage is under varying current), and the
temperature coefficient of $V_Z$ (which can be positive, negative, or near zero depending
on $V_Z$) \cite{sedra2020}. Common power ratings for small-signal Zener diodes are
400\,mW and 1\,W.

\textbf{Purpose of Use.}\quad
The Zener diode's most fundamental use is voltage regulation---maintaining a stable
output voltage from a supply that may vary in magnitude. When connected in reverse bias
across a load with a series current-limiting resistor, the Zener clamps the output voltage
to $V_Z$ regardless of moderate variations in the input supply or load current.

\textbf{Applications.}\quad
Zener diodes provide voltage references in regulated power supplies, protect circuits
from overvoltage transients (clamping), set threshold voltages in comparator circuits,
and stabilize bias voltages in analog amplifier stages \cite{boylestad2015}. They are
found in smartphones, digital multimeters, surge-protected power strips, and battery
management systems.

% -----------------------------------------------------------
\subsection{Light-Emitting Diode (LED)}
\markright{Light-Emitting Diode (LED)}
% -----------------------------------------------------------

\textbf{Definition.}\quad
A light-emitting diode (LED) is a semiconductor p--n junction device that emits light
when forward-biased current flows through it \cite{sedra2020}. As electrons from the
n-type region recombine with holes in the p-type region, they release energy in the
form of photons. The wavelength---and hence color---of the emitted light is determined
by the bandgap energy of the semiconductor material used. LEDs are the most visible
member of the diode family and are used both as indicators and as high-efficiency light
sources.

\imgplaceholder{LED structure diagram showing anode (long lead) and cathode (short lead)}{led}

\textbf{Specifications.}\quad
Key LED specifications include the forward voltage $V_F$ (typically 1.8--2.2\,V for red
or yellow, 3.0--3.5\,V for blue and white), maximum continuous forward current $I_F$
(commonly 20\,mA for standard 5\,mm indicator LEDs), luminous intensity (measured in
millicandela, mcd), viewing angle (typically 15$^\circ$--60$^\circ$), and dominant
wavelength. High-power LEDs for lighting applications carry current ratings of 350\,mA
to several amperes and are mounted on heatsinks \cite{horowitz2015}.

\textbf{Purpose of Use.}\quad
In the laboratory context, LEDs are most commonly used as circuit-state indicators---visual
confirmation that a circuit or a particular branch is energized. They must always be paired
with a current-limiting resistor ($R = (V_\text{supply} - V_F)/I_F$) to prevent
destructive over-current.

\textbf{Applications.}\quad
LEDs are ubiquitous in electronic displays, indicator lights, backlighting for LCD
screens, automotive lighting (headlamps and tail lights), traffic signals, optical fiber
communications (infrared LEDs), remote controls, and solid-state general-purpose lighting.
Their energy efficiency---far superior to incandescent bulbs---has made them the
dominant light source in modern electronics \cite{razavi2013}.

% -----------------------------------------------------------
\subsection{Veroboard}
\markright{Veroboard}
% -----------------------------------------------------------

\textbf{Definition.}\quad
A veroboard (also called a stripboard) is a type of electronics prototyping board
constructed from a single-sided, copper-clad, fiberglass-reinforced epoxy laminate
(FR4) with a regular matrix of holes \cite{wiki:breadboard}. Unlike a breadboard,
a veroboard is intended for semi-permanent circuits: components are soldered in place,
and the continuous copper tracks (strips) running between rows of holes provide electrical
connections. Where connections need to be broken, a track-cutter tool is used to isolate
individual sections.

\imgplaceholder{Veroboard (stripboard) showing copper tracks on the underside}{veroboard}

\textbf{Specifications.}\quad
Standard veroboard has a hole pitch of 2.54\,mm (0.1\,inch), compatible with DIP
integrated circuit packages. The copper strip width is typically 1.5\,mm with 0.5\,mm
gaps between adjacent strips. Board thickness is usually 1.6\,mm (FR4) and the copper
foil weight is 35\,$\mu$m (1\,oz). Unlike breadboards, veroboards can safely handle
continuous currents of several amperes through the copper strips and are suitable for
audio-frequency and low RF circuits where the soldered connections provide reliable
electrical contact \cite{analog:breadboard}.

\textbf{Purpose of Use.}\quad
Veroboard bridges the gap between breadboard prototyping and custom PCB manufacture.
Once a circuit has been verified on a breadboard, it can be transferred to a veroboard
for a more robust, semi-permanent realization. Component placement is planned on paper
or with CAD tools, taking care to plan the track cuts needed to isolate nodes.

\textbf{Applications.}\quad
Veroboards are used in electronics hobbyist projects, small-batch production of simple
circuits, repair workshops (where custom PCBs are not economical), audio amplifier
construction, and educational laboratories where students need a permanent record of
their circuit work \cite{boylestad2015}.

% -----------------------------------------------------------
\subsection{Oscilloscope}
\markright{Oscilloscope}
% -----------------------------------------------------------

\textbf{Definition.}\quad
An oscilloscope is an electronic test instrument that graphically displays the variation
of a voltage signal over time \cite{wiki:oscilloscope}. It plots voltage on the vertical
(Y) axis against time on the horizontal (X) axis, revealing the shape---and therefore
the key parameters---of an electrical signal. Modern digital storage oscilloscopes (DSOs)
sample the analog input, digitize it with an analog-to-digital converter, and store the
waveform in memory for display and analysis.

\imgplaceholder{Digital storage oscilloscope displaying a sinusoidal waveform}{oscilloscope}

\textbf{Specifications.}\quad
The most critical oscilloscope specification is its \textit{bandwidth}---the frequency
at which a sinusoidal input signal is attenuated by 3\,dB (reduced to 70.7\% of its true
amplitude) \cite{tek:oscilloscope}. Laboratory teaching oscilloscopes commonly have
bandwidths of 50--200\,MHz; professional engineering instruments may span 1\,GHz to
over 100\,GHz. The \textit{sample rate} (in samples per second, S/s) must be at least
twice the highest signal frequency (Nyquist criterion); modern DSOs offer 1\,GS/s to
over 100\,GS/s. Additional specifications include memory depth (number of samples stored
per acquisition), vertical sensitivity (typically 1\,mV/div to 10\,V/div), number of
input channels, and trigger options \cite{keysight:bandwidth}.

\textbf{Purpose of Use.}\quad
The oscilloscope allows the engineer to observe signal waveforms that a multimeter cannot
capture, such as rapid transients, oscillation, noise, and timing relationships between
two or more signals. It is essential for verifying the shape of filter outputs, the rise
time of digital pulses, the duty cycle of PWM signals, and the correct timing of serial
communication protocols.

\textbf{Applications.}\quad
Oscilloscopes are used in electronics design and debug labs, manufacturing test, radio
and telecommunications engineering, power electronics (measuring switching transients in
inverters and converters), automotive electronics (decoding CAN bus messages), and
biomedical instrumentation (displaying ECG, EEG, and other biopotential signals)
\cite{wiki:oscilloscope}.

% -----------------------------------------------------------
\subsection{Operational Amplifier (Op-Amp)}
\markright{Operational Amplifier (Op-Amp)}
% -----------------------------------------------------------

\textbf{Definition.}\quad
An operational amplifier (op-amp) is a high-gain, DC-coupled differential voltage
amplifier packaged as an integrated circuit \cite{wiki:opamp}. It has two inputs---an
inverting input ($-$) and a non-inverting input ($+$)---and one output. Its \textit{open-loop}
gain (the gain without any external feedback network) is very large: ideally infinite, and
practically in the range of 20{,}000 to 200{,}000. In virtually all practical applications,
negative feedback is applied via external resistors or other components to set a precise,
stable closed-loop gain.

\imgplaceholder{LM741 op-amp IC (DIP-8 package) with pin diagram and inverting amplifier circuit}{opamp}

\textbf{Specifications.}\quad
For an ideal op-amp: open-loop gain $A_{OL} = \infty$, input impedance $Z_{in} = \infty$,
output impedance $Z_{out} = 0$, bandwidth $= \infty$, and input offset voltage $= 0$
\cite{electronics_tuts:opamp}. A practical general-purpose op-amp such as the LM741
exhibits $A_{OL} \approx 200{,}000$, $Z_{in} \approx 2\,\text{M}\Omega$, $Z_{out} \approx
75\,\Omega$, a unity-gain bandwidth of about 1\,MHz, and an input offset voltage of
1--6\,mV. Other key parameters include the slew rate (maximum rate of output voltage
change, e.g.\ $0.5\,\text{V}/\mu\text{s}$ for the LM741), common-mode rejection ratio
(CMRR, typically 70--90\,dB), and supply voltage range \cite{horowitz2015}.

\textbf{Purpose of Use.}\quad
Op-amps are used to amplify weak signals to useful levels, to implement mathematical
operations (summation, integration, differentiation---whence the name \textit{operational}
amplifier), to compare two voltages, to convert current to voltage, and to buffer a
high-impedance source to drive a low-impedance load without loading the source.

\textbf{Applications.}\quad
Op-amps are found in audio preamplifiers and equalizers, active filter circuits (low-pass,
high-pass, bandpass), instrumentation amplifiers for sensor interfaces (thermocouples,
load cells, strain gauges), analog-to-digital converter front-ends, signal generators
(Wien bridge oscillators, function generators), and feedback control systems in industrial
automation and robotics \cite{sedra2020}.

% -----------------------------------------------------------
\subsection{Bipolar Junction Transistor (BJT)}
\markright{Bipolar Junction Transistor (BJT)}
% -----------------------------------------------------------

\textbf{Definition.}\quad
A bipolar junction transistor (BJT) is a three-terminal semiconductor device that uses
both electrons and holes as charge carriers---hence the term \textit{bipolar}
\cite{boylestad2015}. It consists of two back-to-back p--n junctions sharing a thin common
region called the \textit{base}, flanked by the \textit{emitter} and \textit{collector}
regions. BJTs are manufactured in two complementary configurations: NPN (n-type emitter
and collector surrounding a p-type base) and PNP (p-type emitter and collector surrounding
an n-type base). A small current injected into the base controls a much larger
collector--emitter current, making the BJT a current-controlled device.

\imgplaceholder{NPN BJT schematic symbol, terminal labels, and TO-92 package}{bjt}

\textbf{Specifications.}\quad
The primary BJT specification is the DC current gain $h_{FE}$ (or $\beta$), the ratio
of collector current to base current, typically ranging from 20 to several hundred
depending on device type and operating conditions. Other key parameters include the
collector--emitter breakdown voltage ($V_{CEO}$), the collector--emitter saturation
voltage ($V_{CE(\text{sat})}$ for switching applications), the maximum collector current
$I_C$, the maximum power dissipation $P_C$, and the transition frequency $f_T$
(the frequency at which $h_{FE}$ drops to unity) \cite{sedra2020}.

\textbf{Purpose of Use.}\quad
BJTs are used in two fundamental operating modes. In the \textit{active region}
(base--emitter junction forward biased, base--collector junction reverse biased), the BJT
acts as a linear current amplifier. In the \textit{saturation/cutoff region} (both
junctions forward or reverse biased, respectively), it acts as an electronic switch.
The choice between NPN and PNP depends on the polarity of the supply and the circuit
topology required.

\textbf{Applications.}\quad
BJTs are used in audio amplifiers (common-emitter and push-pull output stages), radio
frequency amplifiers, oscillators, logic circuits (TTL gates), motor drivers, relay
drivers, temperature sensors (the $V_{BE}$ of a BJT is predictably temperature-dependent),
and power regulators \cite{razavi2013}. In modern mixed-signal integrated circuits, BJTs
coexist with CMOS transistors in BiCMOS processes to exploit the high-speed characteristics
of bipolar devices alongside the low power of CMOS.

% -----------------------------------------------------------
\subsection{Metal-Oxide-Semiconductor Field-Effect Transistor (MOSFET)}
\markright{MOSFET}
% -----------------------------------------------------------

\textbf{Definition.}\quad
A MOSFET (metal-oxide-semiconductor field-effect transistor) is a three-terminal
semiconductor device in which the current between the \textit{drain} and \textit{source}
terminals is controlled by the voltage applied to the insulated \textit{gate} terminal
\cite{wiki:mosfet}. The gate is separated from the semiconductor channel by a thin layer
of silicon dioxide ($\text{SiO}_2$), making it electrically insulated; consequently,
virtually no gate current flows, giving the MOSFET a very high input impedance. MOSFETs
are fabricated as n-channel (NMOS) or p-channel (PMOS) types; complementary pairs (CMOS)
are the basis of modern digital logic.

\imgplaceholder{N-channel enhancement MOSFET in TO-220 package with schematic symbol and pin layout}{mosfet}

\textbf{Specifications.}\quad
Key MOSFET parameters include the gate threshold voltage $V_{GS(\text{th})}$ (the
minimum gate-to-source voltage required to form a conductive channel), the drain-to-source
on-state resistance $R_{DS(\text{on})}$ (a few milliohms for power MOSFETs, directly
affecting conduction losses), maximum drain current $I_D$, drain-to-source breakdown
voltage $V_{DSS}$, and transconductance $g_m$ \cite{razavi2013}. For power switching
applications, switching speed (rise/fall times of a few nanoseconds to tens of
nanoseconds) and the gate charge $Q_g$ (determining how quickly the gate can be driven)
are also critical.

\textbf{Purpose of Use.}\quad
In the saturation region, the MOSFET acts as a voltage-controlled current source
suitable for linear amplification. In the cut-off and triode (ohmic) regions, it acts
as an electronic switch---the dominant use in digital and power electronics. Because
MOSFETs switch 50 to 100 times faster than comparable BJTs, they are the preferred
choice in high-frequency power converters \cite{wiki:mosfet}.

\textbf{Applications.}\quad
MOSFETs are the most widely manufactured semiconductor devices in the world. In digital
integrated circuits, billions of CMOS MOSFETs constitute the logic gates, memory cells,
and processors in every computer and smartphone. In power electronics, discrete power
MOSFETs switch thousands of watts in switch-mode power supplies, DC motor drives, solar
inverters, and battery management systems. RF power MOSFETs amplify signals in wireless
base stations, and audio MOSFETs drive high-fidelity loudspeakers \cite{wiki:mosfet}.

% -----------------------------------------------------------
\subsection{Soldering Iron}
\markright{Soldering Iron}
% -----------------------------------------------------------

\textbf{Definition.}\quad
A soldering iron is a hand-held tool with an electrically heated metal tip, called the
\textit{bit}, used to melt solder and join metal parts, wires, or electronic component
leads to circuit boards \cite{wiki:solderiron}. Solder is a low-melting-point metal
alloy (traditionally 60/40 tin--lead, melting at approximately 183\,$^\circ$C; increasingly
lead-free alloys such as Sn--Ag--Cu melting at 217--227\,$^\circ$C) that, when
melted and allowed to solidify, forms a mechanically strong and electrically conductive
joint.

\imgplaceholder{Electrically heated soldering iron with replaceable bit and insulated handle}{solderiron}

\textbf{Specifications.}\quad
Soldering irons are designed to operate within a temperature range of approximately
200\,$^\circ$C to 480\,$^\circ$C (392--896\,$^\circ$F) \cite{wiki:solderiron}. For
general electronics work with lead-free solder, a tip temperature of 315--370\,$^\circ$C
is typical; lower temperatures (300\,$^\circ$C or below) are used for heat-sensitive
SMD components. Power ratings range from about 15\,W for lightweight precision irons to
60--100\,W for heavy-duty or soldering station units. Temperature-controlled soldering
stations include a thermocouple or Curie-point sensor at the tip and closed-loop control
circuitry to maintain the set temperature within a few degrees \cite{wiki:solderiron}.

\textbf{Purpose of Use.}\quad
The soldering iron's purpose is to create permanent, reliable electrical connections by
reflowing solder onto component leads and PCB pads. Good solder joints exhibit a smooth,
shiny appearance and low contact resistance; a poorly made \textit{cold joint} is dull
and grainy, with unreliable conductivity. The soldering iron is also used to desolder
components for rework, using either a desoldering pump or wick.

\textbf{Applications.}\quad
Soldering irons are used in electronics assembly and repair workshops, PCB prototyping
on veroboard, rework and modification of existing PCBs, installation of connectors and
terminal blocks, hobbyist projects, and educational laboratories. High-volume production
lines replace manual soldering with automated wave or reflow soldering ovens, but a
soldering iron remains essential for rework and repair tasks that cannot be automated
\cite{wiki:solderiron}.

% =============================================================
\section{Procedure}
\markright{Procedure}
% =============================================================

\begin{enumerate}
  \item Observe each electronic component and instrument demonstrated in the laboratory.
  \item Carefully examine each component by hand, noting its physical characteristics
        (size, number of terminals, marking, color bands, package type).
  \item Note down the IC or device part number where applicable.
  \item Understand the working principle and important features of each component as
        described in the preceding theory sections.
\end{enumerate}

% =============================================================
\section{Discussion}
\markright{Discussion}
% =============================================================

\vspace{12cm}
% Leave blank for manual entry.

% =============================================================
\section{Conclusion}
\markright{Conclusion}
% =============================================================

\vspace{12cm}
% Leave blank for manual entry.

% =============================================================
% Bibliography
% =============================================================
\markright{References}
\bibliography{references}

% =============================================================
\end{document}
% =============================================================
